\documentclass[11pt,a4paper]{article}
\usepackage[T1]{fontenc}
\usepackage[utf8]{inputenc}
\usepackage{lmodern}
\usepackage[a4paper,top=26mm,bottom=27mm,left=28mm,right=28mm]{geometry}
\usepackage{amsmath,amssymb,amsthm,mathtools}
\usepackage{microtype,booktabs,array,enumitem,xurl}
\usepackage[hidelinks]{hyperref}
\usepackage{fancyhdr}
\hypersetup{pdftitle={Unbounded Integrality Gaps and Additive Hardness for the Skiving Stock Problem},pdfauthor={Eric Winnington},pdfsubject={Combinatorial optimization; bin covering; configuration linear programming}}
\pagestyle{fancy}
\fancyhf{}
\fancyhead[L]{\small E. Winnington}
\fancyhead[R]{\small Integrality gaps for skiving stock}
\fancyfoot[C]{\thepage}
\setlength{\headheight}{14pt}
\renewcommand{\headrulewidth}{0.3pt}
\setlength{\parskip}{3pt}
\setlength{\parindent}{1.2em}
\setlist[enumerate]{topsep=4pt,itemsep=3pt,leftmargin=2em}
\setlist[itemize]{topsep=4pt,itemsep=3pt,leftmargin=1.5em}
\allowdisplaybreaks[1]
\emergencystretch=1.5em
\newtheorem{theorem}{Theorem}
\newtheorem{lemma}[theorem]{Lemma}
\newtheorem{proposition}[theorem]{Proposition}
\newtheorem{corollary}[theorem]{Corollary}
\theoremstyle{definition}
\newtheorem{definition}[theorem]{Definition}
\theoremstyle{remark}
\newtheorem{remark}[theorem]{Remark}
\newtheorem*{sourcegap}{Source Theorem A}
\newtheorem*{sourcehard}{Source Theorem B}
\newcommand{\N}{\mathbb N}
\newcommand{\Z}{\mathbb Z}
\newcommand{\Q}{\mathbb Q}
\newcommand{\R}{\mathbb R}
\newcommand{\BP}{\mathrm{BP}}
\newcommand{\SSP}{\mathrm{SSP}}
\newcommand{\OPT}{\operatorname{OPT}}
\newcommand{\LP}{\operatorname{LP}}
\newcommand{\supp}{\operatorname{supp}}
\newcommand{\card}[1]{\lvert #1\rvert}
\newcommand{\ceil}[1]{\left\lceil #1\right\rceil}
\newcommand{\floor}[1]{\left\lfloor #1\right\rfloor}
\newcommand{\cC}{\mathcal C}
\newcommand{\cU}{\mathcal U}
\newcommand{\cH}{\mathcal H}
\newcommand{\one}{\mathbf 1}
\newcommand{\Pclass}{\mathsf P}
\newcommand{\NPclass}{\mathsf{NP}}
\newcommand{\sourcecommit}{\texttt{adc7f1241b42e322a6451854ab7e4b4c146bf78a}}
\title{\bfseries Unbounded Integrality Gaps and Additive Hardness\\[3pt] for the Skiving Stock Problem}
\author{Eric Winnington}
\date{7 October 2026}

\begin{document}
\maketitle
\thispagestyle{plain}
\begin{abstract}
We prove that the additive integrality gaps of both the ordinary and the proper pattern linear programming relaxations of the one-dimensional skiving stock problem are unbounded. In particular, this refutes Zak's modified integer round-down conjecture for the ordinary relaxation and also disproves its analogue for the stronger proper relaxation. For every fixed integer $c\ge0$, distinguishing instances that cover $B$ bins from instances that cannot cover $B-c$ bins is NP-hard. Consequently, a deterministic polynomial-time algorithm with an absolute additive guarantee exists if and only if $\Pclass=\NPclass$.

The proofs apply a cardinality-saturated complement transfer to the bin-packing construction of OpenAI (24 September 2026). On an instance with $5B$ items, the map $b_i=2-a_i$, with covering threshold $9$, identifies feasible five-item packing and covering configurations. Saturation of the packing LP forces exact use of every item and preserves both covering LP values. Pooling all exceptional groups and unused items converts a covering deficit $d$ into a packing excess at most $\lceil3d/2\rceil$ for the source's size range. More generally, we prove an abstract transfer theorem and exact formulas for both integer optima in terms of a common uniform-hypergraph matching deficiency. The counterexamples and hardness instances can be restricted to an arbitrarily narrow fixed neighborhood of one fifth of the covering threshold. The source-independent transfer is proved here; the source result and the scope of its supplied Lean formalization are identified separately.

This proof was generated by a combination of Anthropic Fable 5.1 and OpenAI GPT 5.6 Sol Pro.
\end{abstract}
\medskip
\noindent\textbf{Keywords:} skiving stock; bin covering; configuration LP; proper relaxation; integrality gap; additive approximation; hypergraph matching.

\section{Introduction}\label{sec:intro}
The one-dimensional skiving stock problem (SSP) asks for the maximum number of disjoint groups of items whose total sizes reach a prescribed threshold. Items may be left unused. The same mathematical problem is usually called \emph{bin covering} in the algorithms literature. Assmann, Johnson, Kleitman, and Leung~\cite{AJKL84} studied this covering counterpart of bin packing in 1984; Johnson, Rennick, and Zak~\cite{JRZ97} subsequently considered skiving in industrial paper processing. Exact algorithms and branch-and-price formulations were developed by Labb\'e, Laporte, and Martello~\cite{LLM95} and Peeters and Degraeve~\cite{PD06}. We use ``bin covering'' and ``skiving stock'' for the single-threshold, offline problem defined in Section~\ref{sec:models}, without additional cardinality or production constraints.

The standard pattern formulation is a maximization integer program with inventory upper bounds. Its continuous relaxation gives an upper bound on the number of covered bins. In analogy with rounding questions for the Gilmore--Gomory cutting-stock formulation~\cite{GG61,ST95}, Zak~\cite{Zak03} proposed the \emph{modified integer round-down property} (MIRDP): the difference between the ordinary LP value and the integer optimum should always be strictly less than two. The proper relaxation disallows a pattern from using more copies of a numerical size than occur in the input, and is a distinct, stronger relaxation~\cite{MS16b}.

\subsection{Main results}
Write $\LP_{\SSP}$ and $\LP_{\SSP,\mathrm{prop}}$ for the ordinary and proper pattern-LP values, respectively, and $\OPT_{\SSP}$ for the maximum number of covered bins. All algorithmic statements use explicitly listed rational items in binary encoding.

\begin{theorem}[Unbounded ordinary and proper gaps]\label{thm:main-gap}
For every integer $c\ge0$, there are a rational SSP instance $J$ and an integer $B$ such that $J$ has exactly $5B$ items and
\begin{equation}\label{eq:main-gap}
 \LP_{\SSP}(J)=\LP_{\SSP,\mathrm{prop}}(J)=B,
 \qquad \OPT_{\SSP}(J)\le B-c-1.
\end{equation}
The covering threshold may be set to $1$, with every item size in
\begin{equation}\label{eq:main-band}
 \left(\frac{53}{270},\frac{11}{54}\right).
\end{equation}
In particular, every covered bin contains at least five items.
\end{theorem}
Taking $c=1$ gives a gap of at least two with an integral LP value. Thus Theorem~\ref{thm:main-gap} refutes Zak's MIRDP for the ordinary relaxation and, separately, the corresponding modified rounding assertion for the proper relaxation. No separation of complexity classes is assumed.

\begin{theorem}[Constant-additive promise hardness]\label{thm:main-hard}
For every fixed integer $c\ge0$, the following promise problem is NP-hard. The input is a rational SSP instance $J$ and an integer $B$, with exactly $5B$ items, threshold $1$, and all sizes in~\eqref{eq:main-band}. Distinguish
\[
 \textnormal{YES: }\OPT_{\SSP}(J)=B,
 \qquad
 \textnormal{NO: }\OPT_{\SSP}(J)\le B-c-1.
\]
\end{theorem}

\begin{corollary}[Absolute additive approximation]\label{cor:main-complexity}
There exist an absolute constant $C\ge0$ and a deterministic polynomial-time algorithm returning a feasible filling with at least $\OPT_{\SSP}(J)-C$ covered bins on every explicitly encoded rational instance $J$ if and only if $\Pclass=\NPclass$. Under this equality, an optimal filling can be constructed in polynomial time.
\end{corollary}

The proofs use two results from OpenAI's bin-packing manuscript~\cite{OAI26}: unbounded configuration-LP gaps with integral LP value, and NP-hardness of every fixed additive packing allowance. The abstract transfer is developed first, independently of that source. Section~\ref{sec:source} then records the precise source interface and proves Theorems~\ref{thm:main-gap} and~\ref{thm:main-hard}. The stronger source-specific bound $a_i<7/30$ is needed for the improved conversion factor and the displayed size band, but not for the existence of an unbounded-gap transfer.

\subsection{Previous bounds and related work}\label{sec:history}
Martinovic and Scheithauer studied ordinary integer rounding~\cite{MS16a}, the proper relaxation~\cite{MS16b}, and alternative integer formulations~\cite{MS16c}. Their later work treated divisible instances~\cite{MS17,MS18a}, polyhedral sufficient conditions for IRDP~\cite{MS18b}, and hypergraph matching formulations~\cite{MS18c}. Kartak and Ripatti~\cite{KR19} reported ordinary SSP gaps of approximately $1.1795$ and proper SSP gaps of approximately $1.1319$, improving the previous ordinary record of approximately $1.0476$. These are the largest published lower-bound values located in the literature reviewed for this paper; the numbers quoted are the authors' reported rounded values, not new computations of their examples.

For $m$ numerical item types, Martinovic and Scheithauer~\cite{MS19} proved the general bound $\Delta<\lceil m/2\rceil$ and sharper bounds depending on residual instances and item lengths. Such bounds do not assert a universal constant independent of $m$. In the divisible case, where the threshold is an integer multiple of every size, the upper bound improved to $3/2$ in 2017~\cite{MS17} and to $4/3$ in Martinovic's paper~\cite{M22}, published online in 2020 and in the 2022 volume of \emph{4OR}. These special-case bounds are not contradicted by our construction.

Packing--covering relationships, including LP-aware relationships, predate this work. Zak's volume-tight LP equivalence~\cite[Theorem~2]{Zak03} is restated in~\cite[Lemma~8]{MS19}: equality of the ordinary packing LP to normalized total volume is equivalent to the corresponding equality for skiving. For integral volume,~\cite[Theorem~7]{MS19} transfers failure of the ordinary integer round-up property to failure of the ordinary integer round-down property, and conversely. Our theorem uses \emph{cardinality saturation} and affine complementation instead of volume-tightness; it controls arbitrary additive deficits and explicitly treats proper columns. It does not assume that total original item size equals $B$.

Hypergraph representations of SSP are also established~\cite{MS18c}. The shared matching-deficiency formulas in Section~\ref{sec:matching} identify the additional structure used here: every $(k-1)$-item subset packs, every $(k+1)$-item subset covers, and the two systems have the same feasible $k$-item subsets. Related cardinality-constrained covering models~\cite{EIL13} impose different feasibility restrictions. No claim of novelty is made for the arithmetic identity $\sum_i(s-a_i)=s|M|-\sum_i a_i$ itself. The contribution is its precise LP-saturated transfer, its quantitative consequences, and its application to the cited 2026 construction. We did not locate a previous statement combining these features.

Subsequent computational work includes improved flow formulations~\cite{MDISS20}, stabilized column generation combined with the Reflect Arc-Flow model~\cite{KIMS23}, and variable-sized cutting and skiving formulations~\cite{DM26}. These developments concern formulations and solution methods rather than establishing the unbounded ordinary and proper gaps proved here. The asymptotic approximation schemes of Csirik, Johnson, and Kenyon~\cite{CJK01} and Jansen and Solis-Oba~\cite{JSO03} remain compatible with Theorem~\ref{thm:main-hard}: a guarantee with a nonzero multiplicative loss is not an absolute additive guarantee.

\section{Models and rounding conventions}\label{sec:models}
An item list has physical index set $V=\{1,\ldots,N\}$. Distinct indices represent distinct copies even when their numerical sizes coincide. We assume $N>0$ when discussing size minima; the empty instance has integer and LP values zero. For a type vector $v$, write $|v|=\sum_\sigma v_\sigma$. All comparisons of vectors are componentwise.

\subsection{Packing configurations}
A bin-packing instance $I$ consists of $a_i\in\Q\cap(0,1]$. Let $S$ be its set of numerical sizes and $n_\sigma\ge1$ their multiplicities. The unrestricted configuration family is
\[
 \cC(I)=\left\{v\in\Z_{\ge0}^{S}:\sum_{\sigma\in S}\sigma v_\sigma\le1\right\}.
\]
It is finite and downward closed. Its members may exceed the input multiplicity of a type. Define
\begin{equation}\label{eq:packing-lp}
 \LP_{\BP}(I)=\min\left\{\sum_{v\in\cC(I)}z_v:
 z\ge0,\quad \sum_{v\in\cC(I)}v_\sigma z_v\ge n_\sigma\ (\sigma\in S)\right\}.
\end{equation}
The proper configuration family is $\cC_{\mathrm{prop}}(I)=\{v\in\cC(I):v\le n\}$. The individual-copy family is
\[
 \cH_{\BP}(I)=\left\{H\subseteq V:\sum_{i\in H}a_i\le1\right\}.
\]
Its LP minimizes $\sum_H z_H$ subject to $z\ge0$ and $\sum_{H\ni i}z_H\ge1$ for every $i$. Its value is $\LP_{\BP,\mathrm{ind}}(I)$. Physical packings partition all items; their minimum size is $\OPT_{\BP}(I)$.

\subsection{Covering patterns and minimality}
For an SSP instance $J$, let $b_i\in\Q_{>0}$ and threshold $L\in\Q_{>0}$. With numerical type set $T$ and multiplicities $n_\tau$, define the upward-closed family
\[
 \cU(J)=\left\{p\in\Z_{\ge0}^{T}:\sum_{\tau\in T}\tau p_\tau\ge L\right\}.
\]
A pattern is \emph{minimal} if no strictly smaller nonnegative integer vector belongs to $\cU(J)$. Write $\cU_{\min}(J)$ for these patterns and $\cU_{\min,\mathrm{prop}}(J)=\{p\in\cU_{\min}(J):p\le n\}$. The ordinary relaxation is
\begin{equation}\label{eq:covering-lp}
 \LP_{\SSP}(J)=\max\left\{\sum_{p\in\cU_{\min}(J)}y_p:
 y\ge0,\quad \sum_{p\in\cU_{\min}(J)}p_\tau y_p\le n_\tau\ (\tau\in T)\right\}.
\end{equation}
Restricting the columns to $\cU_{\min,\mathrm{prop}}(J)$ defines $\LP_{\SSP,\mathrm{prop}}(J)$. The integer version equals $\OPT_{\SSP}(J)$: every selected pattern consumes inventory, and unused items need not be assigned.

\begin{lemma}[Finite minimal-pattern formulation]\label{lem:minimal}
Restricting covering columns to minimal patterns changes neither the integer optimum nor either LP optimum. The minimal pattern family is finite.
\end{lemma}
\begin{proof}
Every feasible pattern contains a minimal subpattern, obtained by deleting items while coverage remains feasible. Replacing a column by such a subpattern preserves its objective contribution and decreases resource consumption. Properness is preserved. For a minimal nonzero pattern $p$, choose a type $\tau$ with $p_\tau>0$. Minimality gives $b\cdot p<L+\tau\le L+b_{\max}$. Hence $|p|<(L+b_{\max})/b_{\min}$, proving finiteness. The replacement argument applies to every finite-support solution of the unrestricted formulation. It also applies to nonnegative summable solutions by aggregating their weights onto the finite minimal family; bounded inventory bounds total weight. Conversely, every minimal-pattern solution is feasible in the unrestricted family.
\end{proof}
This formulation and equivalence agree with the standard minimal-pattern convention; see~\cite[Section~2]{MS19}. In particular, the presence of infinitely many dominated nonminimal columns is not a different relaxation.

\begin{lemma}[Proper types and individual copies]\label{lem:lifting}
For packing, the proper type LP equals the individual-copy LP. For covering, the proper type LP equals the LP on feasible subsets of physical copies, with constraints $\sum_{H\ni i}y_H\le1$.
\end{lemma}
\begin{proof}
Projecting each physical subset to its type counts and adding coincident column weights preserves feasibility and objective value. Conversely, for a proper type column $p$, independently choose a uniformly random $p_\sigma$-subset of the $n_\sigma$ copies of every type $\sigma$, and distribute the column's weight over all these choices. A fixed copy of type $\sigma$ is selected with probability $p_\sigma/n_\sigma$. Its total coverage or consumption is therefore the type total divided by $n_\sigma$, which satisfies the corresponding inequality. Every generated subset is feasible because feasibility depends only on type counts. Nonminimal covering subsets can subsequently be reduced by Lemma~\ref{lem:minimal}. This is an existence argument; it does not require efficient enumeration of all lifted columns.
\end{proof}
Consequently,
\[
 \LP_{\BP}\le\LP_{\BP,\mathrm{ind}}\le\OPT_{\BP},
 \qquad
 \OPT_{\SSP}\le\LP_{\SSP,\mathrm{prop}}\le\LP_{\SSP}.
\]

\subsection{The modified round-down assertion}
The ordinary SSP gap is $\Delta(J)=\LP_{\SSP}(J)-\OPT_{\SSP}(J)$; the proper gap replaces $\LP_{\SSP}$ by $\LP_{\SSP,\mathrm{prop}}$. MIRDP means
\begin{equation}\label{eq:mirdp}
 \Delta(J)<2,
 \quad\text{equivalently}\quad
 \OPT_{\SSP}(J)\ge\floor{\LP_{\SSP}(J)}-1.
\end{equation}
Zak's conjecture concerns the ordinary relaxation~\cite{Zak03,MS16a}. We refer separately to the \emph{proper analogue} of~\eqref{eq:mirdp}; we do not identify the two formulations historically or mathematically. When both LP values equal an integer $B$, an optimum at most $B-2$ violates both assertions.

\section{A cardinality-saturated transfer theorem}\label{sec:transfer}
We first isolate the combinatorial hypotheses, without using numerical item sizes or any particular source construction.

\begin{definition}[Complementary configuration systems]\label{def:systems}
Fix a finite type set $S$, an inventory $n\in\Z_{>0}^{S}$ with $|n|=kB$, and an integer $k\ge2$. Let $\cC\subset\Z_{\ge0}^{S}$ be a finite downward-closed packing family containing every unit vector, and let $\cU\subseteq\Z_{\ge0}^{S}$ be an upward-closed covering family with a finite minimal family. Assume
\begin{equation}\label{eq:cardinality-hypotheses}
 |v|\le k\ (v\in\cC),\qquad |p|\ge k\ (p\in\cU),
 \qquad
 \{v\in\cC:|v|=k\}=\{p\in\cU:|p|=k\}.
\end{equation}
Physical subsets are feasible when their type-count vectors belong to the corresponding family. Define the ordinary and proper LPs as in Section~\ref{sec:models}, and denote the two integer optima by $\OPT_{\cC}$ and $\OPT_{\cU}$.
\end{definition}
The bounds in~\eqref{eq:cardinality-hypotheses} concern \emph{unrestricted type columns}, not just subsets of available copies. This distinction is necessary for conclusions about the ordinary relaxation.

\begin{theorem}[LP saturation transfers in both directions]\label{thm:lp-transfer}
Under Definition~\ref{def:systems},
\begin{align}
 \LP_{\cC}=B&\quad\Longleftrightarrow\quad\LP_{\cU}=B,\label{eq:ordinary-equiv}\\
 \LP_{\cC,\mathrm{prop}}=B&\quad\Longleftrightarrow\quad\LP_{\cU,\mathrm{prop}}=B.\label{eq:proper-equiv}
\end{align}
In particular, $\LP_{\cC,\mathrm{ind}}=B$ implies that all four LP values in~\eqref{eq:ordinary-equiv}--\eqref{eq:proper-equiv} equal $B$.
\end{theorem}
\begin{proof}
Summing inventory constraints and using cardinality gives
\begin{equation}\label{eq:universal-lp-bounds}
 \LP_{\cC}\ge B,\qquad \LP_{\cU}\le B,
\end{equation}
and the same inequalities hold for the proper families. The packing programs attain their optima because their column families are finite and a bounded objective sublevel is compact; the covering programs are bounded finite LPs on minimal columns.

Let $z$ be a packing solution of value $B$. Then
\begin{equation}\label{eq:saturation}
 kB=\sum_\sigma n_\sigma
 \le\sum_v|v|z_v
 \le k\sum_vz_v=kB.
\end{equation}
Every positive column consequently has cardinality $k$, and every type constraint holds with equality. Indeed, the two differences in~\eqref{eq:saturation} are sums of nonnegative terms. The common $k$-item layer in~\eqref{eq:cardinality-hypotheses} therefore turns $z$ into a covering solution of the same value with exact inventory use. Its columns are minimal covering patterns, because no smaller pattern covers. Together with~\eqref{eq:universal-lp-bounds}, this proves the forward implication in~\eqref{eq:ordinary-equiv}.

Conversely, a covering solution $y$ of value $B$ satisfies
\[
 kB=k\sum_py_p\le\sum_p|p|y_p\le\sum_\sigma n_\sigma=kB.
\]
Every positive column has cardinality $k$ and every resource constraint is tight. The same columns form a packing solution of value $B$, proving the reverse implication. Both arguments preserve properness, proving~\eqref{eq:proper-equiv}. Finally, Lemma~\ref{lem:lifting}, the inclusions of proper families, and~\eqref{eq:universal-lp-bounds} yield the last assertion.
\end{proof}

\begin{remark}[Physical-copy proof and repetitions]\label{rem:copies}
When $\LP_{\cC,\mathrm{ind}}=B$, the corresponding equality is
\[
 kB\le\sum_{i\in V}\sum_{H\ni i}z_H
 =\sum_H|H|z_H\le k\sum_Hz_H=kB.
\]
Thus each positive subset has $k$ items and every individual-copy constraint binds. Aggregation gives exact type consumption, regardless of repeated numerical sizes. This argument neither assumes distinct sizes nor turns an improper type column into a proper one. Empty columns have zero weight at equality. Nonminimal covering columns cannot improve the cardinality upper bound.
\end{remark}

\begin{theorem}[Constructive deficit conversion]\label{thm:conversion}
Under Definition~\ref{def:systems}, suppose every physical subset with at most $r$ items packs, where $1\le r\le k$. If a covering solution has $q=B-d$ groups, of which $g$ contain more than $k$ items, then
\begin{equation}\label{eq:general-conversion}
 \OPT_{\cC}\le B+\ceil{\frac{(k-r)(g+d)}r}
 \le B+\ceil{\frac{(k-r)(k+1)}r\,d}.
\end{equation}
Conversely, every packing with $B+e$ bins supplies a covering with at least $B-(k-1)e$ groups. These conversions operate in polynomial time on explicitly listed solutions and do not require any LP hypothesis.
\end{theorem}
\begin{proof}
Keep the $q-g$ covering groups with exactly $k$ items as packing bins. Pool \emph{all} other items: those in the $g$ larger groups and those unused by the covering. The pool size is exactly
\begin{equation}\label{eq:pool-size}
 kB-k(q-g)=k(g+d).
\end{equation}
The larger groups use at least $(k+1)g$ pooled items. Consequently
\begin{equation}\label{eq:g-bound}
 (k+1)g\le k(g+d),\qquad g\le kd.
\end{equation}
Partition the pool arbitrarily into sets of at most $r$ items. These sets pack by hypothesis. The total bin count is at most
\[
 B-d-g+\ceil{\frac{k(g+d)}r}
 =B+\ceil{\frac{(k-r)(g+d)}r},
\]
which proves~\eqref{eq:general-conversion} using~\eqref{eq:g-bound}.

For the converse, let $f$ packing bins have exactly $k$ items. Every other bin has at most $k-1$, so
\[
 kB\le kf+(k-1)(B+e-f)=(k-1)(B+e)+f.
\]
Therefore $f\ge B-(k-1)e$. Keep these $f$ groups as covering groups and leave all remaining items unused. Grouping or discarding explicitly represented items takes polynomial time. No test for an optimal packing or covering is involved.
\end{proof}

\begin{corollary}[Two-sided gap comparison on the same pairs]\label{cor:gap-equiv}
Set $e=\OPT_{\cC}-B$ and $d=B-\OPT_{\cU}$. Under Definition~\ref{def:systems},
\begin{equation}\label{eq:two-sided}
 0\le d\le(k-1)e,
 \qquad 0\le e\le(k^2-1)d.
\end{equation}
For any family of such pairs with fixed $k$ and the appropriate LP values equal to $B$, additive packing gaps are unbounded if and only if additive covering gaps are unbounded. This assertion applies to the ordinary values, or to the proper values, when their respective saturation hypothesis holds.
\end{corollary}
\begin{proof}
Use Theorem~\ref{thm:conversion} with optimal solutions and $r=1$; singletons pack by Definition~\ref{def:systems}. Cardinality implies $e,d\ge0$. Under saturation, $e$ and $d$ are exactly the corresponding additive gaps.
\end{proof}
The equivalence is restricted to these complementary pairs. It does not assert that an arbitrary packing instance and an arbitrary skiving instance have comparable gaps.

\subsection{Affine complements realize the hypotheses}
\begin{proposition}[Scalar complement and admissible shifts]\label{prop:complement}
Let $I$ have $kB$ rational sizes $a_i\in(0,1]$, and suppose every unrestricted packing configuration has at most $k$ items. Write $\alpha=\min_i a_i$ and choose a fixed rational $s$ satisfying
\begin{equation}\label{eq:shift}
 s>\max_i a_i,\qquad s+(k-1)\alpha>1.
\end{equation}
Define $J_s$ by $b_i=s-a_i$ and $L_s=ks-1$. Then $L_s>0$, all complemented sizes are positive, and the two configuration systems satisfy Definition~\ref{def:systems}. In particular, $s=2$ always suffices. The map preserves numerical type multiplicities and has polynomial bit complexity.
\end{proposition}
\begin{proof}
For any multiset $M$ of cardinality $m$, including repeated type entries,
\begin{equation}\label{eq:complement-identity}
 \sum_{i\in M}(s-a_i)\ge ks-1
 \quad\Longleftrightarrow\quad
 \sum_{i\in M}a_i\le1-(k-m)s.
\end{equation}
The largest complemented size is $s-\alpha$. By~\eqref{eq:shift},
\[
 (k-1)(s-\alpha)<ks-1.
\]
Thus no multiset of fewer than $k$ items covers. For $m=k$,~\eqref{eq:complement-identity} is exactly the unit-capacity packing condition. The other family properties follow from the packing and covering definitions and Lemma~\ref{lem:minimal}. Positivity of $L_s$ follows from $ks>s+(k-1)\alpha>1$. The map on sizes is injective. For fixed rational $s$, subtraction and multiplication by the fixed normalizing factor have polynomial bit cost. For $s=2$, a fraction $u/v$ becomes $(2v-u)/v$.
\end{proof}
A convenient numerical condition for the packing-cardinality hypothesis is $a_i>1/(k+1)$ for every numerical size: it rules out $k+1$ items even in unrestricted columns. By contrast, positivity of $s-a_i$ alone is not enough for the covering-cardinality conclusion. For example, five items of size $17/100$ and $s=21/100$ give complemented size $1/25$ and threshold $1/20$, already covered by two items. Condition~\eqref{eq:shift} excludes this failure.

\section{Exact optima from a shared matching deficiency}\label{sec:matching}
The conversion becomes particularly transparent when all immediately smaller packing groups and all immediately larger covering groups are automatically feasible.

\begin{theorem}[Exact matching-deficiency formulas]\label{thm:matching}
Under Definition~\ref{def:systems}, assume every physical subset of at most $k-1$ items packs and every physical subset of exactly $k+1$ items covers. Let $\cH$ be the $k$-uniform hypergraph on $V$ whose edges are the common feasible $k$-item subsets. Write $\nu=\nu(\cH)$ for its maximum matching cardinality and $h=B-\nu$. Then
\begin{equation}\label{eq:exact-optima}
 \OPT_{\cC}=B+\ceil{\frac h{k-1}},
 \qquad
 \OPT_{\cU}=B-\ceil{\frac h{k+1}}.
\end{equation}
No LP assumption is required for these identities.
\end{theorem}
\begin{proof}
Let a packing use $P$ bins, of which $f$ contain $k$ items. These $f$ subsets form a matching, so $f\le\nu$. Hence
\[
 kB\le kf+(k-1)(P-f)\le(k-1)P+\nu,
\]
and $P\ge\lceil(kB-\nu)/(k-1)\rceil$. Conversely, use a maximum matching and partition the remaining items into sets of at most $k-1$. This attains
\[
 \nu+\ceil{\frac{k(B-\nu)}{k-1}}
 =B+\ceil{\frac h{k-1}}.
\]

For a covering with $Q$ groups, let $f$ have cardinality $k$. Again $f\le\nu$. Each other group uses at least $k+1$ items, so
\[
 kB\ge kf+(k+1)(Q-f)\ge(k+1)Q-\nu.
\]
Therefore $Q\le\lfloor(kB+\nu)/(k+1)\rfloor$. A maximum matching followed by arbitrary disjoint $(k+1)$-item groups from the remainder attains
\[
 \nu+\floor{\frac{k(B-\nu)}{k+1}}
 =B-\ceil{\frac h{k+1}}.
\]
Any final remainder is left unused. These constructions prove both equalities.
\end{proof}

A fractional perfect matching in $\cH$ is a nonnegative assignment to its edges giving total incidence one at every physical item. The saturation proof shows that such a matching exists if and only if the proper packing LP equals $B$, equivalently the proper covering LP equals $B$. When saturation holds,~\eqref{eq:exact-optima} identifies the same integral matching deficiency behind both additive gaps. This is a statement about the full feasible hypergraph, not only the support of a selected fractional solution.

\begin{corollary}[The five-item window]\label{cor:five}
Let $I$ contain $5B$ items with $1/6<a_i<7/30$, and let $J$ have sizes $2-a_i$ and threshold $9$. Every covering solution with $B-d$ groups, of which $g$ have more than five items, yields a packing using at most
\begin{equation}\label{eq:three-halves}
 B+\ceil{\frac{g+d}{4}}
 \le B+\ceil{\frac{3d}{2}}.
\end{equation}
Moreover, with $h=B-\nu(\cH)$ for the feasible five-item hypergraph,
\begin{equation}\label{eq:five-exact}
 \OPT_{\BP}(I)-B=\ceil{h/4},
 \qquad B-\OPT_{\SSP}(J)=\ceil{h/6}.
\end{equation}
\end{corollary}
\begin{proof}
Every four original items have total size below $28/30<1$, while every six complemented items have total size greater than $6(2-7/30)>9$. Apply Theorems~\ref{thm:conversion} and~\ref{thm:matching} with $k=5$ and $r=4$. In the direct conversion, the pool contains $5(g+d)$ items, so its bin count is
\[
 B-d-g+\ceil{\frac{5(g+d)}4}=B+\ceil{\frac{g+d}4}.
\]
The inequality $g\le5d$ yields~\eqref{eq:three-halves}.
\end{proof}
Whenever $h$ is unbounded, the ratio of the two quantities in~\eqref{eq:five-exact} tends to $3/2$. Thus no coefficient strictly smaller than $3/2$, even with a fixed additive correction, can bound packing excess by covering deficit along every unbounded-gap family in this window. This sharpness concerns optimal values; it does not claim that finding a maximum matching is easy.

\section{The source results and the main applications}\label{sec:source}
We use the manuscript \emph{Additive hardness and unbounded configuration gaps in bin packing}, dated 24 September 2026~\cite{OAI26}. All repository references in this paper are pinned to commit
\begin{center}\sourcecommit.\end{center}
The source's proof and formalization are external inputs, not results independently re-established here. The elementary transfer above does not use the source's trees, interval gadgets, role labels, or extraction argument.

\subsection{The exact interface}
\begin{sourcegap}[\cite{OAI26}, unbounded configuration-gap theorem]
For every integer $t\ge0$, there are a rational bin-packing instance $I$ and an integer $B$ such that
\begin{equation}\label{eq:source-gap}
 N=5B,\qquad \frac16<a_i<1,\qquad
 \LP_{\BP}(I)=\LP_{\BP,\mathrm{ind}}(I)=B,
 \qquad \OPT_{\BP}(I)>B+t.
\end{equation}
\end{sourcegap}

\begin{sourcehard}[\cite{OAI26}, constant-additive hardness theorem]
For every fixed integer $t\ge0$, it is NP-hard to distinguish a packing into $B$ bins from the absence of a packing into $B+t$ bins. The reduction outputs an explicitly listed rational instance with $N=5B$ and $a_i>1/6$; every item also satisfies $a_i\le1$. The output and running time are polynomial in the input's binary encoding length for each fixed $t$.
\end{sourcehard}
No assertion about LP values on the hardness instances is needed or assumed. In particular, Source Theorem~B is not obtained by applying Source Theorem~A as an algorithmic black box.

The source's graph-to-packing construction is parameterized by a fixed allowance $t$ and a fixed tolerance $\rho>0$. A graph with a vertex cover of size at most $K$ yields a $B$-bin packing; a packing into $B+t$ bins yields a cover of size at most $K+\rho n$. The source applies this to a \emph{gap} vertex-cover promise obtained from a satisfiability gap. An arbitrary no-instance of the exact vertex-cover decision problem, with optimum merely greater than $K$, would not suffice. Its unconditional LP-gap construction instead uses the four-vertex clique and a fractional witness. These are separate applications of the construction~\cite[Sections~5--7]{OAI26}.

\subsection{The stronger size bound}
Both source families may be chosen from the actual scalar construction, which satisfies
\begin{equation}\label{eq:strong-source-window}
 \frac16<a_i<\frac7{30}.
\end{equation}
For completeness, the estimate used here is explicit. In the source's notation,
\[
 a_i=\frac15+\gamma p(i)+\mu q(i)+\eta w(i),
 \qquad
 \gamma=\frac1{1000P_0},\quad
 \mu=\frac{\gamma}{1000Q_n},\quad
 \eta=\frac{\mu}{1000},
\]
where $P_0=100^{11}$, $Q_n\ge1$, $|p(i)|<P_0$, $|q(i)|\le2Q_n$, and $|w(i)|\le6$. We have renamed the source's last scalar coefficient to $\eta$ to avoid confusion with the matching number $\nu$. Therefore
\begin{equation}\label{eq:source-estimate}
 \left|a_i-\frac15\right|
 \le\frac1{1000}+\frac{2}{10^6P_0}
       +\frac{6}{10^9P_0Q_n}
 <\frac1{30}.
\end{equation}
This is the scalar-encoding estimate in~\cite[Section~3]{OAI26}. It proves~\eqref{eq:strong-source-window} for the constructions supplying both source results. It is not an extra premise silently appended to their exported statements.

\subsection{Formalization scope and dependency boundary}\label{sec:formal-boundary}
The source repository provides a Lean scope document and a comparator challenge for the theorem \texttt{OAI.BinPackingGap.main\_results}~\cite{OAILean26}. The exported structural statement includes $N=5B$, $a_i>1/6$, both LP equalities, and the strict packing lower bound. The exported hardness structure includes the item count, the lower size bound, and a polynomial-time binary-encoded promise reduction. The comparator configuration points to the separate solution module \texttt{OAI.Computability.}\allowbreak\texttt{BinPacking.Main}. The placeholder proof in the challenge file is not the solution proof.

The upper bound $7/30$ belongs to the construction's size analysis, rather than the exported comparator statement. It is used only in Corollary~\ref{cor:five}, the improved choice of source allowance, and the size-window refinements. The unbounded-gap and hardness transfer itself already follows from the weaker interface using Theorem~\ref{thm:conversion} with $r=1$, giving allowance $24c$.

This paper does not claim an independent rebuild of the source Lean development or an independently executed comparator check. The transfer proved here is also not Lean-formalized. Appendix~\ref{app:source} lists the exact source locations so that these two verification tasks can be separated.

\subsection{Proofs of the main theorems}
\begin{proof}[Proof of Theorem~\ref{thm:main-gap}]
Fix $c\ge0$ and set
\begin{equation}\label{eq:source-allowance}
 t=\ceil{\frac{3c}{2}}.
\end{equation}
Choose the instance of Source Theorem~A from the construction satisfying~\eqref{eq:strong-source-window}. Form the complement $J$ with $b_i=2-a_i$ and threshold $9$. Every unrestricted packing configuration has at most five items. Proposition~\ref{prop:complement} and Theorem~\ref{thm:lp-transfer} give
\[
 \LP_{\SSP}(J)=\LP_{\SSP,\mathrm{prop}}(J)=B.
\]
Suppose a covering uses $q\ge B-c$ groups. Cardinality implies $q\le B$, so its deficit $d=B-q$ satisfies $0\le d\le c$. Corollary~\ref{cor:five} gives
\[
 \OPT_{\BP}(I)\le B+\ceil{\frac{3d}{2}}
 \le B+\ceil{\frac{3c}{2}}=B+t,
\]
contradicting Source Theorem~A. Thus $\OPT_{\SSP}(J)\le B-c-1$. This argument does not require discarding groups to a potentially negative target $B-c$.

Finally, divide all complemented sizes and the threshold by $9$. From~\eqref{eq:strong-source-window},
\[
 \frac{2-a_i}{9}\in\left(\frac{53}{270},\frac{11}{54}\right).
\]
Scaling preserves every feasible pattern and all integer and LP values. This proves all assertions.
\end{proof}
For $c=1$, the sufficient source condition is $\OPT_{\BP}>B+2$, that is, an integer packing gap of at least three. A source gap of exactly two is not enough for this conclusion from~\eqref{eq:three-halves}.

\begin{proof}[Proof of Theorem~\ref{thm:main-hard}]
Fix $c$ and use Source Theorem~B with the fixed parameter~\eqref{eq:source-allowance}, choosing the source construction with~\eqref{eq:strong-source-window}. Compose its reduction with complementation and normalization.

In the yes-case, a $B$-bin packing of $5B$ items, with at most five per bin, has exactly five in each bin. Each becomes a covering group, and no covering can have more than $B$ groups. Thus $\OPT_{\SSP}=B$.

In the no-case, a covering with $q\ge B-c$ groups would have deficit $d\le c$ and would yield a packing with at most $B+\lceil3c/2\rceil$ bins by Corollary~\ref{cor:five}. This contradicts source soundness. Hence $\OPT_{\SSP}\le B-c-1$.

For fixed $c$, the source parameter $t$ is fixed. All construction constants, including the source's finite trees and any tolerance chosen in its hardness proof, are constants of the reduction rather than runtime input. The source outputs polynomially many rational items with polynomial total bit length. The complement map adds only polynomial bit work, proving the required polynomial-time reduction. No LP is solved in this proof.
\end{proof}
The result concerns rational, binary-encoded data. Clearing denominators preserves the mathematical instance and has polynomial bit cost, but may produce very large integer magnitudes; we make no separate claim of strong NP-hardness for polynomially bounded integer lengths.

\begin{proof}[Proof of Corollary~\ref{cor:main-complexity}]
Suppose the asserted algorithm and constant $C$ exist. Fix an integer $c\ge C$. In the yes-case of Theorem~\ref{thm:main-hard}, its output has at least $B-C\ge B-c$ groups. In the no-case, no feasible output has that many. Counting the output groups decides an NP-hard promise problem in polynomial time, implying $\Pclass=\NPclass$.

Conversely, suppose $\Pclass=\NPclass$. For a rational instance with $N$ items and target $q\le N$, a certificate assigns every item a label in $\{0,1,\ldots,q\}$, where $0$ means unused. Exact rational arithmetic verifies that each label $1,\ldots,q$ has total size at least the threshold. The certificate and the verification have polynomial bit complexity. The decision problem therefore belongs to $\NPclass$. The same is true when a prefix of item assignments is prescribed.

Binary search over $0\le q\le N$ finds the optimum using polynomial-time decisions. To construct a filling of that value, process the items successively. For each item, try every label in $\{0,1,\ldots,q\}$ and retain one for which a feasible completion of the current prefix exists. At least one such label exists by the invariant. At most $N(N+1)$ completion tests produce an optimal filling. Handle the empty instance directly. This is a deterministic polynomial-time constructive algorithm with $C=0$.
\end{proof}

\section{Consequences and limitations}\label{sec:discussion}
\subsection{Arbitrarily narrow fixed size windows}
The band in Theorem~\ref{thm:main-gap} is not an essential width restriction.
\begin{corollary}[Concentration near one fifth]\label{cor:narrow}
For every fixed rational $\varepsilon>0$, Theorems~\ref{thm:main-gap} and~\ref{thm:main-hard} remain valid with all normalized SSP item sizes in
\[
 (1/5-\varepsilon,\,1/5+\varepsilon)\cap(0,1).
\]
The source allowance may still be chosen as $\lceil3c/2\rceil$.
\end{corollary}
\begin{proof}
Use a fixed rational shift $s\ge2$ and normalize by $5s-1$. Proposition~\ref{prop:complement} applies. The normalized deviation is
\[
 \frac{s-a_i}{5s-1}-\frac15
 =\frac{1-5a_i}{5(5s-1)}.
\]
By~\eqref{eq:strong-source-window}, $|1-5a_i|<1/6$, giving a deviation below $1/[30(5s-1)]$. Choose the fixed $s$ large enough that this is below $\varepsilon$. Original four-item subsets still pack, so Theorem~\ref{thm:conversion} gives the same $3/2$ factor. The chosen normalized sizes are in $(0,1)$, and all additional constants depend only on the fixed window.
\end{proof}
The statement is about a \emph{fixed} window. It neither places a uniform bound on denominators nor supplies small explicit numerical counterexamples.

\subsection{Growth rates and special cases}
Theorems~\ref{thm:main-gap} and~\ref{thm:main-hard} give unbounded additive phenomena, but no asymptotically growing lower bound expressed in terms of the number of items or encoding length. This is also a limitation stated for the source construction~\cite{OAI26}. On the complementary five-item subclass, a logarithmic upper bound does follow: Hoberg and Rothvoss~\cite{HR17} bound packing excess above the configuration LP by $O(\log(2+B))$, and~\eqref{eq:five-exact} gives
\[
 B-\OPT_{\SSP}(J)=\ceil{h/6}\le\ceil{h/4}
 =\OPT_{\BP}(I)-B.
\]
Thus the ordinary and proper SSP gaps on the saturated subclass considered here are $O(\log(2+B))$. This does not establish a logarithmic upper bound for arbitrary SSP instances, nor a matching logarithmic lower bound.

The divisible-case $4/3$ bound~\cite{M22} remains intact because the complement does not preserve divisibility. The fixed-type bound $\Delta<\lceil m/2\rceil$~\cite{MS19} likewise shows that unbounded gaps require an unbounded number of numerical sizes. Affine complementation preserves that number exactly. For completeness, Appendix~\ref{app:boundaries} records two boundary observations that distinguish these results from easier one-sided size intervals and the fixed-type regime.

\subsection{What is and is not settled}
The ordinary MIRDP and its proper analogue are disproved by Theorem~\ref{thm:main-gap}; no uniform absolute additive approximation is possible unless $\Pclass=\NPclass$ by Theorem~\ref{thm:main-hard}. These conclusions do not depend on a claimed novelty of the complement identity. The open quantitative issue is the possible rate of gap growth in general SSP instances. Smaller-cardinality analogues would require packing families satisfying the stated cardinality and LP hypotheses; they do not follow merely by replacing five by three or four in the source application.

\paragraph{AI assistance and verification.}
ChatGPT assisted with drafting, literature organization, and analysis of the proofs. Such assistance is not a formal verification. The proof status and source dependence are described in Section~\ref{sec:formal-boundary}; the source's Lean development was not independently rebuilt in preparing this manuscript.

\appendix
\section{Source locators and logical dependencies}\label{app:source}
The following paths are relative to the pinned \texttt{openai/math} repository. They identify the external premises rather than substituting for verification of their proofs.
\begin{center}
\small
\begin{tabular}{@{}>{\raggedright\arraybackslash}p{0.27\textwidth}>{\raggedright\arraybackslash}p{0.37\textwidth}>{\raggedright\arraybackslash}p{0.28\textwidth}@{}}
\toprule
Premise & Source location & Role in this paper\\
\midrule
$5B$ items; both packing LP values $B$; arbitrary strict additive gap &
\texttt{lean/OAI/Computability/}\newline\texttt{BinPacking/Main.lean}; first conjunct of \texttt{main\_results} &
Ordinary and proper unbounded-gap application\\[5pt]
Fixed-allowance promise hardness; $5B$ items; $a_i>1/6$ &
\texttt{lean/ComparatorChallenges/}\newline\texttt{BinPackingGap.lean}; \texttt{PackingGapReduction} and \texttt{PackingGapNPHard} &
Hardness application; no LP equality needed\\[5pt]
Actual sizes below $7/30$ &
Source paper, Section~3, scalar-encoding proof; \nolinkurl{build/source/instance.tex} in the preprint directory &
Four-item pooling, exact matching formulas, narrow bands\\[5pt]
Formal theorem and solution mapping &
\texttt{lean/ComparatorChallenges/}\newline\texttt{BinPackingGap.json}; \nolinkurl{lean/docs/118.md} &
Scope and comparator metadata only\\
\bottomrule
\end{tabular}
\end{center}
The configuration definitions in the comparator distinguish unrestricted numerical-type vectors from subsets of physical indices. The former are capacity-bounded, not inventory-bounded. The source's fractional witness constructs the latter first and then aggregates numerical coincidences. The cardinality bound gives both matching lower bounds. This is why repeated numerical sizes do not create a convention mismatch in the present use of Source Theorem~A.

The solution module is \texttt{OAI.Computability.BinPacking.Main}. The configuration lists \texttt{propext}, \texttt{Quot.sound}, and \texttt{Classical.choice} as permitted axioms. The \texttt{sorry} placeholder in the challenge specifies the target to be proved; it is not the proof in the named solution module. These observations describe the repository layout, not a report of a completed local comparator run.

\section{Boundary cases}\label{app:boundaries}
\subsection{One-sided reciprocal intervals have IRDP}
\begin{proposition}\label{prop:one-sided}
Let $r\ge1$ be an integer. If all SSP item sizes belong to $(L/(r+1),L/r]$, the ordinary and proper relaxations both have the integer round-down property.
\end{proposition}
\begin{proof}
Let $t$ items have size exactly $L/r$, and let $u$ have smaller size. Fewer than $r$ items never cover. A group of $r$ items covers only when all of them have size $L/r$, while every group of $r+1$ items covers. The ordinary LP value is
\[
 F=\frac tr+\frac{u}{r+1}.
\]
Indeed, pure-type columns of cardinality $r$ for the largest size and $r+1$ for every smaller size give this value fractionally. For the matching dual upper bound, assign weight $1/r$ to the largest size and $1/(r+1)$ to every smaller size. Every feasible covering pattern has weight at least one.

Write $t=rq+v$ with $0\le v<r$. Cover $q$ bins using largest-size items and group the remaining $v+u$ items $r+1$ at a time. The number obtained is
\[
 q+\floor{\frac{v+u}{r+1}}=\floor{F}.
\]
For the equality, the difference between $v/r+u/(r+1)$ and $(v+u)/(r+1)$ is $v/[r(r+1)]<1/(r+1)$; the latter has denominator $r+1$, so their floors agree. The LP bound proves optimality. The proper LP lies between this integer value and $F$, so its floor is the same integer optimum.
\end{proof}
Thus hard reciprocal-size windows must be distinguished from these one-sided intervals. In particular, intervals $(L/4,L/3]$ and $(L/5,L/4]$ are not consequences of the present hardness theorem.

\subsection{Fixed types and high multiplicity}
The general fixed-dimensional integer-cone theorem of Goemans and Rothvoss~\cite{GR20} also applies to SSP with a fixed number $m$ of sizes and binary-encoded multiplicities. Here is the direct formulation. For a target $T$, consider the rational polytopes in dimension $m+1$
\[
 \begin{aligned}
 P&=\{(x,1):0\le x\le n,\ b^{\mathsf T}x\ge L\},\\
 Q_T&=\{(u,t):0\le u\le n,\ T\le t\le N\},
 \qquad N=\sum_jn_j.
 \end{aligned}
\]
An integer point of $P$ is a proper covering pattern with one group counted in the last coordinate. A sum of integer points of $P$ lies in $Q_T$ exactly when at least $T$ groups can be covered within the inventory. The fixed-dimensional theorem decides this intersection and supplies a compact integer-cone representation. Binary search over $T\in\{0,\ldots,N\}$ uses polynomially many tests in the binary input length. The compact representation is essential in the high-multiplicity model: an expanded list of all physical copies need not have polynomial length. This observation is separate from the explicit-list complexity convention used in the main theorems.

\section{A sharper broad-window conversion for five items}\label{app:nineteen}
Without the upper bound $a_i<7/30$, Theorem~\ref{thm:conversion} gives coefficient $24$ for $k=5$. Scalar feasibility improves this to $19$.
\begin{proposition}\label{prop:nineteen}
Let $I$ contain $5B$ items with $a_i\in(1/6,1]$, and complement by $b_i=2-a_i$, $L=9$. A covering with $B-d$ groups yields a packing using at most $B+19d$ bins. The coefficient $19$ is sharp under these hypotheses alone; this sharpness assertion does not impose LP saturation.
\end{proposition}
\begin{proof}
Let $g_6$ groups contain six items and $g_{\ge7}$ contain at least seven. The total cardinality surplus above five gives
\[
 g_6+2g_{\ge7}\le5d.
\]
A six-item covering group has original total size at most three. It contains a pair of total size at most one: otherwise any partition of its six items into three pairs would have total size greater than three. Keep every five-item covering group as a packing bin. In the pooled remainder, pack one feasible pair from each six-item group, and pack all other items singly. Starting from the singleton-pool count yields
\[
 \OPT_{\BP}\le B+4d+3g_6+4g_{\ge7}
 \le B+4d+3(g_6+2g_{\ge7})\le B+19d.
\]
A feasible pair can be found by checking the finitely many pairs within each six-item group.

For sharpness, take $25d$ items of size $501/1000$ and $5d$ of size $495/1000$, where $d\ge1$, and put $B=6d$. Five larger items and one smaller item have total size three, so their complements cover exactly threshold nine. These give $5d=B-d$ groups. No five-item group covers, since every five original items have total size greater than one, so this covering is optimal. The $25d$ larger packing items require distinct bins; all smaller items can be paired with larger ones. Hence $\OPT_{\BP}=25d=B+19d$. The packing LP in this example is not $B$, so the example does not constrain improvements obtainable from additional saturation assumptions.
\end{proof}

\bibliographystyle{plain}
\bibliography{references}
\end{document}
